% =====================================================================
%  Chapter 12 — The \idxw{Yarkovsky} Drift
% =====================================================================

\chapter{The Yarkovsky Drift}
\label{ch:yarkovsky}

A body that is heated unevenly and cooled unevenly does not receive and
re-emit sunlight symmetrically, and it therefore experiences a small,
persistent thrust. The effect is called Yarkovsky\idx{Yarkovsky} after the scientist who
identified it in the nineteenth century, and it is the reason the orbits of
small asteroids are not closed ellipses.

For a body of this class the effect is not a curiosity. It is the
accumulated along-track displacement that decides whether a keyhole is hit.

\section{The mechanism}

The body's rotation carries surface elements through the day and night
cycle. A surface element that is heated and then cools radiates at
temperatures that depend on when it was last illuminated, and the radiation
leaves preferentially in the direction normal to the surface. The net recoil
over a rotation is a small force along the spin axis and a small torque
about it.

The force depends on the obliquity, on the thermal properties of the surface
and on the spin rate. Writing it in the form used in the 2012 study of
Yarkovsky\idx{Yarkovsky} and YORP,
\begin{equation}
\frac{d a}{d t} \propto
\frac{(1-A)\cos\gamma}{\sqrt{a\,\rho_b\,D}}\,
f(\Theta, \omega),
\label{eq:yarkovsky}
\end{equation}
where \(A\) is the albedo, \(\gamma\) is the angle between the spin axis and
the ecliptic normal, \(\rho_b\) is the bulk density, \(D\) is the diameter and
\(\omega\) the rotation rate.

Three of the five terms are physical properties that a characterisation
campaign measures: the albedo, the bulk density and the diameter. The
fourth, the spin axis, comes from the lightcurve. The fifth, the thermal
parameter \(\Theta\), does not.

\section{The thermal parameter, and why the book will not extrapolate it}

The function \(f(\Theta, \omega)\) contains the thermal parameter\idx{Thermal parameter}
\begin{equation}
\Theta = \frac{\Gamma\,\omega}{\varepsilon\,\sigma\,T_{\mathrm{SSP}}^{3}},
\label{eq:theta}
\end{equation}
where \(\Gamma\) is a shape-and-latitude factor of order unity,
\(\varepsilon\) the emissivity, \(\sigma\) the Stefan-Boltzmann constant and
\(T_{\mathrm{SSP}}\) the temperature of a surface element at the
steady-state temperature of the sub-solar point. For a body with a low
thermal inertia, \(T_{\mathrm{SSP}}\propto D^{-1/2}\), so \(\Theta\) grows
steeply as the body shrinks.

This is the reason the book refuses to extrapolate an absolute drift rate
across the size range between the calibration object and the class of
interest. The scaling laws in \eqref{eq:yarkovsky} isolate the explicit size
and density factor, and that factor can be quoted as a ratio. The thermal
parameter cannot be scaled that way, because it varies with size through a
different and stronger dependence.

\begin{deeper}{The rule, and why it exists}
The 2012 calibration is the asteroid pseudo-Geographos: 2,560 metres in
diameter, a bulk density of 2,500 kilogrammes per cubic metre, a semimajor
axis of 1.246 astronomical units, a Yarkovsky\idx{Yarkovsky} drift of 18.1 metres per year
and a YORP acceleration of \(2.26\times10^{-3}\) radians per year squared.
Applying the scaling laws to a 60-metre body gives ratios of the order of
forty to a hundred for the drift and of several hundred to several thousand
for the torque. Those ratios are the size-and-density part of the effect and
nothing else. The published absolute drift rates for real objects are used
wherever an absolute drift appears in this book, and no absolute drift is
derived by scaling.
\end{deeper}

\begin{table}[!htbp]
\centering
\footnotesize
\caption{The size-and-density part of the scaling laws, expressed as ratios
to the pseudo-Geographos calibration. The thermal-parameter factor is
excluded, and the ratios are therefore lower bounds on the true scaling.}
\label{tab:rg-ratios}
\begin{tabularx}{\textwidth}{@{}P{0.99in}P{0.86in}P{0.86in}Y@{}}
\toprule
\rowcolor{panel}
\textbf{Diameter and density} & \textbf{Yarkovsky} & \textbf{YORP} &
\textbf{Comment} \\
\midrule
51 m, 1,300 & \(\times 96\) & \(\times 4{,}827\) & the YR4 fast-rotating
model with a rubble pile \\
\rowcolor{panel}
51 m, 2,800 & \(\times 45\) & \(\times 2{,}241\) & as above, coherent \\
60.8 m, 1,300 & \(\times 81\) & \(\times 3{,}409\) & the YR4
thermophysical model, rubble pile \\
\rowcolor{panel}
60.8 m, 2,800 & \(\times 38\) & \(\times 1{,}583\) & as above, coherent \\
140 m, 1,300 & \(\times 35\) & \(\times 643\) & the catalogue's threshold
class \\
\rowcolor{panel}
140 m, 2,800 & \(\times 16\) & \(\times 299\) & as above, coherent \\
\bottomrule
\end{tabularx}
\end{table}

The ratios in \tabref{tab:rg-ratios} are large because the calibration object
is two orders of magnitude larger, and \eqref{eq:yarkovsky} makes the
dependence explicit. The practical consequence is that the smaller the body,
the more the drift rate depends on knowing its size and density --- which is
precisely the class where the catalogue knows neither.

\section{What a shape measurement does, and does not, buy here}

This is the point at which an over-statement is easiest to make, and the
published evidence is unusually clear.

The same 2012 study finds that gross shape differences move the Yarkovsky
drift by only about ten per cent for objects of the same size. What enhances
the drift appreciably is \emph{roughness beaming} --- the effect of a rough
surface on the thermal emission --- which can raise it by tens of per cent
and, in the most favourable cases, by a factor of two. And it finds that the
sensitivity to the size is three times the sensitivity to the shape: a three
per cent error in the diameter produces a three per cent error in Yarkovsky,
while the same three per cent in the shape produces only about one per cent
in the drift.

The conclusion is uncomfortable but it is the right one, and the book states
it rather than avoiding it: \textbf{a resolved shape measurement does not
substantially improve the Yarkovsky drift}. It improves the YORP torque,
which Chapter 13 takes up, and it supplies the volume on which the entire
accounting chain rests, and it supplies the spin pole that enters
\eqref{eq:yarkovsky} through \(\cos\gamma\). But the drift itself is
governed by the size, the density and the thermal parameter.

This is why the argument of the book is anchored on the accounting chain and
on the close encounter, and not on the drift. The drift is where the shape
measurement earns its keep through the spin state and the density, and not
through the shape's effect on the thermal emission.

\section{What the live case says}

The 2025 JWST analysis of 2024~YR4 makes the connection explicit in its own
discussion. The object's seasonal heating is negligible for the March 2025
observations, because of the near-equator-on viewing geometry; the seasonal
Yarkovsky drift is expected to dominate the diurnal component and might
become measurable during future encounters; and the non-gravitational forces
such as the Yarkovsky effect contribute to the growing uncertainty in the
future trajectory of the object, making very long-term predictions
impossible.

The wording matters. The drift is not yet measurable for this object. It is
expected to become measurable at a future encounter. Until then it is a
term in a covariance matrix that grows with the arc, and it is one of the
reasons the JWST analysis explicitly declines to offer a long-term impact
probability.

\begin{keytakeaways}
\item The Yarkovsky drift depends on the albedo, the obliquity, the density,
the diameter and a thermal parameter\idx{Thermal parameter} that does not scale with them.
\item The scaling laws give the size-and-density ratios, which for this class
are between sixteen and ninety-six times the calibration object, but the
book does not extrapolate an absolute drift across that range.
\item A resolved shape measurement improves the drift by only about ten per
cent; what it improves materially is YORP, the volume and the pole.
\item For 2024 YR4 the drift is not yet measurable and is expected to become
measurable at a future encounter, and it is already a growing term in the
trajectory uncertainty.
\end{keytakeaways}
